\section{Risultati}

\textbf{Classificazione.} Il circuito verificato ha profondità 412, cioè 412 strati
di porte in sequenza. Contiene 224 porte CX, che invertono un qubit in base al valore di un altro.
Nei due scenari UC1/UC2, le SVM raggiungono F1 $=0{,}919/0{,}923$ e AUC $=0{,}965/0{,}958$.
Su 30 repliche per scenario, tutte le strategie fattorizzano entro il limite previsto.

La sola ricerca TOP-4 porta i cicli medi a $1{,}00$, contro $1{,}37/1{,}50$ del Metodo~1.
Il confronto è significativo ($p_{\mathrm{Holm}}=0{,}0024/0{,}0015$).

Il Metodo~2 richiede $1{,}50/1{,}37$ cicli e non migliora significativamente il Metodo~1
($p_{\mathrm{Holm}}=0{,}7314/0{,}0874$). Richiede più cicli di TOP-4
($p_{\mathrm{Holm}}=0{,}0033/0{,}0096$).
Buone metriche di classificazione non garantiscono quindi una riduzione del costo necessario a trovare i fattori.

La zero-noise extrapolation non migliora il riferimento.
La variante su tre livelli triplica il costo di campionamento.

\textbf{Correzione d'errore.} Il codice a ripetizione a $d=3$ riproduce i valori previsti dalla teoria.
Per Steane, l'errore dopo la correzione cresce circa con il quadrato dell'errore fisico (esponente stimato $1{,}94$).
La pseudo-soglia di Steane è $p\simeq0{,}08$: qui gli errori dopo la
correzione eguagliano quelli del qubit non protetto.
Riguarda un codice a distanza fissata.

Per il surface code, le soglie stimate sono $0{,}86\%$ ($Z$) e $0{,}81\%$ ($X$).
$Z$ e $X$ indicano due modi di preparare e misurare il qubit.

La rete usata da sola non supera il matching nelle quattordici configurazioni iniziali ($d=3,5$).

Con crosstalk (interferenza fra qubit), la rete affiancata al matching riduce $p_L$ del $18$--$21\%$ a $d=3$ e dell'$1$--$3\%$ a $d=5$.
A $d=7$ non emerge un miglioramento significativo.
In questi test, alcune correlazioni non sono rappresentate dal grafo del matching e lasciano un margine per la rete.

\textbf{Sensibilità di Shor.} Il successo senza rumore è $74{,}89\%$, con intervallo al 95\%
$[74{,}59\%;\ 75{,}18\%]$, coerente con il $75\%$ teorico.

Il successo si stabilizza intorno a $p_g=0{,}1$; a $p_g=0{,}5$ vale $24{,}59\%$.
È vicino al pavimento uniforme del $24{,}61\%$, dovuto agli esiti accettati anche senza informazione utile.

\textbf{QFT approssimata.} Su $N=21$ eliminare le piccole rotazioni durante
i calcoli aritmetici riduce le porte a due qubit da $49\,661$ a $23\,178$.
Il successo senza rumore scende però dal $46{,}67\%$ al $40{,}13\%$.
Nelle simulazioni rumorose eseguite non emerge un vantaggio misurabile.

Un modello semplificato stima un guadagno massimo di circa tre punti percentuali.
Richiede errori delle porte a due qubit circa duecento volte inferiori alla calibrazione considerata.

Nella stima di fase, tutti i nove confronti favoriscono la QFT troncata.
In sei confronti, l'intervallo al 95\% del miglioramento è interamente positivo.
Un intervallo interamente positivo distingue il guadagno da zero nel singolo confronto.
Il maggior aumento è di 12,6 punti percentuali su dieci qubit
(intervallo al 95\%: 9,8--15,3 punti).
La scelta migliore conserva le rotazioni con gli angoli più grandi,
limitandoli a due o tre valori nelle dimensioni esaminate.
